\section{Step 4: Response}\label{sec:response}

A per-flow Benign FPR of a few percent (3.8\%) looks modest in a classification table. At a gateway that blocks host pairs, it means blocked innocent hosts. We therefore score response in false blocks: a block is false if the blocked host pair carries no attack flow in that window.

\emph{ARP watcher.} For the attack that flow features cannot see (Section~\ref{sec:splits}), a small watcher raises a high alert when a MAC other than the known one announces the router IP (.1), the classic arpwatch idea~\cite{arpwatch}. In the MITM window it raised one high alert 1.7\,s into the window, naming dc:a6:32:dc:27:d5, and none in the other nine windows. This shows that the rule works here; it is not a false-alarm rate.

\emph{Policies.} With enforcement on, the gateway blocks a host pair (OVS drop rules for both directions) once the pair has collected $k$ attack verdicts; optionally, it ignores MITM-ARP flow verdicts, which the ARP watcher covers. We chose the policy on validation data by simulating the verdict stream: among policies within 0.005 of the best share of attack flows that start after their pair is blocked, we took the one with the fewest false blocks. DDoS-SYN flows, the bulk of all attack flows, dominate this criterion. On validation, $k=1$ gave 128 false blocks, and $k=10$ with MITM-ARP verdicts ignored gave 1, while still covering 0.995 of attack flows (0.996 for $k=1$).

\begin{table}[t]
\centering
\caption{Live blocking over the ten test windows (final harness, blocks cleared per window). Scoring false blocks by the triggering flow gives the same counts.}
\label{tab:blocking}
\footnotesize
\setlength{\tabcolsep}{3pt}
\begin{tabular}{lrrrr}
\hline
Policy & Blocked & Correct & False & Packets dropped \\
\hline
Block on first verdict & 79 & 19 & 60 & 326,268 \\
$k=10$, MITM-ARP ignored & 9 & 9 & 0 & 163,546 \\
\hline
\end{tabular}
\end{table}

\emph{Live results.} Table~\ref{tab:blocking} gives the two enforced runs. Blocking on the first verdict blocked all 7 Mirai bots (about 100k packets, 35\% of the window, dropped) but also 60 benign pairs, 32 of them on a DoS-HTTP and 21 on a MITM-ARP verdict. The false blocks dropped only 132 packets, but each cuts a benign pair off for the rest of the window. The validation-chosen policy blocked all six DDoS attackers (133k packets dropped), the DoS-HTTP attacker (30k), the brute-force attacker (217) and the port scanner (0; the scan had ended), with no false blocks. It missed the 7 Mirai bots and the 3 web attackers, whose pairs produce too few flows to reach 10 verdicts. In the block-on-first-verdict run the DDoS-SYN window was replayed at only 9,045 packets/s (52\%); since slower replay splits flows (Section~\ref{sec:fidelity}), that run's DDoS-window flow and block counts may be inflated.

\begin{table}[t]
\centering
\caption{Trade-off over $k$, simulated per window on the logged run-1 verdict stream (26 attacker pairs). Time to block: median over blocked attacker pairs, from the pair's first flow. Per-class columns: attacker pairs blocked.}
\label{tab:kcurve}
\footnotesize
\setlength{\tabcolsep}{2.2pt}
\begin{tabular}{rcrrcccccr}
\hline
$k$ & MITM ign. & Blocked & False & Attackers & Mirai & Web & BruteF & Time (s) \\
\hline
1  & no  & 83 & 64 & 19/26 & 7/7 & 3/3 & 1/1 & 19.0 \\
1  & yes & 56 & 37 & 19/26 & 7/7 & 3/3 & 1/1 & 19.0 \\
2  & yes & 23 &  5 & 18/26 & 6/7 & 3/3 & 1/1 & 33.8 \\
3  & yes & 14 &  1 & 13/26 & 1/7 & 3/3 & 1/1 & 15.8 \\
5  & yes & 10 &  0 & 10/26 & 0/7 & 1/3 & 1/1 & 15.3 \\
10 & yes &  9 &  0 &  9/26 & 0/7 & 0/3 & 1/1 & 15.3 \\
20 & yes &  8 &  0 &  8/26 & 0/7 & 0/3 & 0/1 & 15.3 \\
\hline
\end{tabular}
\end{table}

\emph{The full trade-off.} To show the whole curve, we replayed the policy logic on the verdicts logged in run~1 (Table~\ref{tab:kcurve}; this run did not enforce, so its $k=1$ counts differ slightly from the enforced run). Ignoring MITM-ARP verdicts alone cuts false blocks at $k=1$ from 64 to 37. With $k=3$, 1 false block remains, and with $k\ge5$ none, but Mirai is lost from $k=3$ on (1/7, then 0/7), web attackers from $k=5$ (1/3, then 0/3 at $k\ge10$), and the brute-force attacker at $k=20$. All six DDoS-SYN pairs are blocked at every $k$, and no MITM-ARP pair at any $k$. The median time to block (15.3--33.8\,s) is not monotonic in $k$, because the set of blocked pairs changes with $k$. It does not grow, because the pairs still blocked at high $k$ are mostly DDoS-SYN pairs (6 of 9 at $k=10$), whose floods quickly yield many verdicts.

\emph{Late blocks.} In the run-1 simulation, at every $k$, only 11.1\% of the attack flows in the live windows start after their pair is blocked; DDoS-SYN flows dominate this share. A block needs a verdict, and a verdict needs a flow to end (idle timeout 15\,s), so most attack flows in a 60\,s window arrive before any block can act. Blocks still cut the later packets of long attack flows: the enforced runs dropped 0.503 ($k=1$) and 0.342 (policy) of the DDoS-SYN window's packets, 0.707 and 0.698 of the DoS-HTTP window's, and 0.347 of the Mirai window's ($k=1$). The validation share (0.995) is far too optimistic: the validation simulation blocks a pair at the last packet of its $k$-th flow predicted as an attack, ignoring the export delay. This is a weakness of our selection criterion; the false-block counts, which decided the choice, do not depend on block timing.
